\chapter{ItemStacks und Inventare}
\label{ch:items/inventories}

In diesem Kapitel lernst du, wie man Inventare\index{inventar@Inventar}
verwendet und verändert -- egal ob Spielerinventar, Node-Inventar oder
freistehendes Inventar.

\section{Was sind ItemStacks und
Inventare?}\label{was-sind-itemstacks-und-inventare}

Ein ItemStack\index{itemstack@\texttt{ItemStack}} sind die Daten hinter
einem einzelnen Feld eines Inventars.

Ein \emph{Inventar} ist eine Sammlung von \emph{Inventarlisten}, von
denen jede ein 2D-Gitter aus ItemStacks ist. Im Zusammenhang mit
Inventaren heißen Inventarlisten einfach \emph{Listen}.

Spieler und Nodes haben jeweils nur ein einziges Inventar; Listen
erlauben mehrere Gitter innerhalb dieses Inventars. Standardmäßig hat
der Spieler die Liste \texttt{main} für den Großteil seines Inventars
und einige Listen für das Crafting.

\section{ItemStacks}\label{itemstacks}

ItemStacks bestehen aus vier Teilen: \texttt{name}, \texttt{count},
\texttt{wear} und Metadaten.\index{metadaten@Metadaten}

Der Itemname kann der Name eines registrierten Items, ein
Alias\index{alias@Alias} oder ein unbekannter Itemname sein. Unbekannte
Items entstehen häufig, wenn Benutzer Mods deinstallieren oder Mods
Items ohne Vorkehrungen wie Aliase entfernen.

\begin{Shaded}
\begin{Highlighting}[]
\FunctionTok{print}\OperatorTok{(}\VariableTok{stack}\OperatorTok{:}\NormalTok{get\_name}\OperatorTok{())}
\VariableTok{stack}\OperatorTok{:}\NormalTok{set\_name}\OperatorTok{(}\StringTok{"default:dirt"}\OperatorTok{)}

\ControlFlowTok{if} \KeywordTok{not} \VariableTok{stack}\OperatorTok{:}\NormalTok{is\_known}\OperatorTok{()} \ControlFlowTok{then}
    \FunctionTok{print}\OperatorTok{(}\StringTok{"Unbekanntes Item!"}\OperatorTok{)}
\ControlFlowTok{end}
\end{Highlighting}
\end{Shaded}

Die Anzahl ist immer 0 oder größer. Im normalen Spielverlauf sollte sie
die maximale Stapelgröße des Items, \texttt{stack\_\allowbreak{}max},
nicht überschreiten. Admin-Befehle und fehlerhafte Mods können aber dazu
führen, dass Stacks größer werden.

\begin{Shaded}
\begin{Highlighting}[]
\FunctionTok{print}\OperatorTok{(}\VariableTok{stack}\OperatorTok{:}\NormalTok{get\_stack\_max}\OperatorTok{())}
\end{Highlighting}
\end{Shaded}

Ein ItemStack kann leer sein; dann ist die Anzahl 0.

\begin{Shaded}
\begin{Highlighting}[]
\FunctionTok{print}\OperatorTok{(}\VariableTok{stack}\OperatorTok{:}\NormalTok{get\_count}\OperatorTok{())}
\VariableTok{stack}\OperatorTok{:}\NormalTok{set\_count}\OperatorTok{(}\DecValTok{10}\OperatorTok{)}
\end{Highlighting}
\end{Shaded}

ItemStacks lassen sich mit der Funktion \texttt{ItemStack} auf
verschiedene Weise erzeugen:

\begin{Shaded}
\begin{Highlighting}[]
\NormalTok{ItemStack}\OperatorTok{()} \CommentTok{{-}{-} name="", count=0}
\NormalTok{ItemStack}\OperatorTok{(}\StringTok{"default:pick\_stone"}\OperatorTok{)} \CommentTok{{-}{-} count=1}
\NormalTok{ItemStack}\OperatorTok{(}\StringTok{"default:stone 30"}\OperatorTok{)}
\NormalTok{ItemStack}\OperatorTok{(\{} \VariableTok{name} \OperatorTok{=} \StringTok{"default:wood"}\OperatorTok{,} \VariableTok{count} \OperatorTok{=} \DecValTok{10} \OperatorTok{\})}
\end{Highlighting}
\end{Shaded}

Item-Metadaten sind ein unbegrenzter Schlüssel-Wert-Speicher für Daten
zum Item. Schlüssel-Wert heißt: Über einen Namen (den Schlüssel,
\emph{key}) greifst du auf die Daten (den Wert, \emph{value}) zu. Manche
Schlüssel haben eine besondere Bedeutung, etwa \texttt{description}, mit
dem ein Stack eine eigene Beschreibung bekommt. Mehr dazu im Kapitel
Speicherung und Metadaten~(Kapitel~\ref{ch:map/storage}).

\section{Inventarorte}\label{inventarorte}

Der Inventarort legt fest, wo und wie ein Inventar gespeichert wird. Es
gibt drei Arten: Spieler, Node und freistehend. Ein Inventar ist fest an
genau einen Ort gebunden -- Änderungen am Inventar wirken sofort.

\subsection{Node-Inventare}\label{node-inventare}

Node-Inventare gehören zur Position eines bestimmten Nodes, etwa einer
Truhe. Der Node muss geladen sein, weil das Inventar in den
Node-Metadaten~(Kapitel~\ref{ch:map/storage}) gespeichert wird.

\begin{Shaded}
\begin{Highlighting}[]
\VariableTok{on\_punch} \OperatorTok{=} \KeywordTok{function}\OperatorTok{(}\VariableTok{pos}\OperatorTok{,} \VariableTok{node}\OperatorTok{)}
    \KeywordTok{local} \VariableTok{inv} \OperatorTok{=}
        \VariableTok{core}\OperatorTok{.}\NormalTok{get\_inventory}\OperatorTok{(\{} \FunctionTok{type} \OperatorTok{=} \StringTok{"node"}\OperatorTok{,} \VariableTok{pos} \OperatorTok{=} \VariableTok{pos} \OperatorTok{\})}
    \CommentTok{{-}{-} jetzt das Inventar verwenden}
\KeywordTok{end}\OperatorTok{,}
\end{Highlighting}
\end{Shaded}

So erhältst du eine \emph{Inventar-Referenz}, meist \emph{InvRef}
genannt. Mit Inventar-Referenzen veränderst du ein Inventar.
\emph{Referenz} bedeutet, dass die Daten nicht im Objekt selbst liegen;
das Objekt ändert die Daten direkt an Ort und Stelle.

Den Ort einer Inventar-Referenz ermittelst du so:

\begin{Shaded}
\begin{Highlighting}[]
\KeywordTok{local} \VariableTok{location} \OperatorTok{=} \VariableTok{inv}\OperatorTok{:}\NormalTok{get\_location}\OperatorTok{()}
\end{Highlighting}
\end{Shaded}

\subsection{Spielerinventare}\label{spielerinventare}

Spielerinventare erhältst du auf ähnliche Weise oder über eine
Spieler-Referenz. Der Spieler muss online sein, damit du auf sein
Inventar zugreifen kannst.

\begin{Shaded}
\begin{Highlighting}[]
\KeywordTok{local} \VariableTok{inv} \OperatorTok{=} \VariableTok{core}\OperatorTok{.}\NormalTok{get\_inventory}\OperatorTok{(\{}
    \FunctionTok{type} \OperatorTok{=} \StringTok{"player"}\OperatorTok{,}
    \VariableTok{name} \OperatorTok{=} \StringTok{"player1"}\OperatorTok{,}
\OperatorTok{\})}
\CommentTok{{-}{-} oder}
\KeywordTok{local} \VariableTok{inv} \OperatorTok{=} \VariableTok{player}\OperatorTok{:}\NormalTok{get\_inventory}\OperatorTok{()}
\end{Highlighting}
\end{Shaded}

\subsection{Freistehende Inventare}\label{freistehende-inventare}

Ein freistehendes Inventar (\emph{detached inventory}) ist unabhängig
von Spielern und Nodes. Es wird bei einem Neustart nicht gespeichert.

\begin{Shaded}
\begin{Highlighting}[]
\KeywordTok{local} \VariableTok{inv} \OperatorTok{=} \VariableTok{core}\OperatorTok{.}\NormalTok{get\_inventory}\OperatorTok{(\{}
    \FunctionTok{type} \OperatorTok{=} \StringTok{"detached"}\OperatorTok{,}
    \VariableTok{name} \OperatorTok{=} \StringTok{"inventory\_name"}\OperatorTok{,}
\OperatorTok{\})}
\end{Highlighting}
\end{Shaded}

Anders als die anderen Inventararten musst du ein freistehendes Inventar
erst anlegen, bevor du darauf zugreifst:

\begin{Shaded}
\begin{Highlighting}[]
\VariableTok{core}\OperatorTok{.}\NormalTok{create\_detached\_inventory}\OperatorTok{(}\StringTok{"inventory\_name"}\OperatorTok{)}
\end{Highlighting}
\end{Shaded}

\texttt{create\_\allowbreak{}detached\_\allowbreak{}inventory} nimmt
drei Argumente entgegen, von denen nur das erste -- der Inventarname --
Pflicht ist. Das zweite ist eine Tabelle mit
Callbacks,\index{callback@Callback} mit denen du steuerst, wie Spieler
mit dem Inventar umgehen können:

\begin{Shaded}
\begin{Highlighting}[]
\CommentTok{{-}{-} Freistehendes Inventar, in das man nur hineinlegen kann}
\VariableTok{core}\OperatorTok{.}\NormalTok{create\_detached\_inventory}\OperatorTok{(}\StringTok{"inventory\_name"}\OperatorTok{,} \OperatorTok{\{}
    \VariableTok{allow\_move} \OperatorTok{=} \KeywordTok{function}\OperatorTok{(}
        \VariableTok{inv}\OperatorTok{,}
        \VariableTok{from\_list}\OperatorTok{,}
        \VariableTok{from\_index}\OperatorTok{,}
        \VariableTok{to\_list}\OperatorTok{,}
        \VariableTok{to\_index}\OperatorTok{,}
        \VariableTok{count}\OperatorTok{,}
        \VariableTok{player}
    \OperatorTok{)}
        \ControlFlowTok{return} \VariableTok{count} \CommentTok{{-}{-} Verschieben erlauben}
    \KeywordTok{end}\OperatorTok{,}

    \VariableTok{allow\_put} \OperatorTok{=} \KeywordTok{function}\OperatorTok{(}
        \VariableTok{inv}\OperatorTok{,}
        \VariableTok{listname}\OperatorTok{,}
        \VariableTok{index}\OperatorTok{,}
        \VariableTok{stack}\OperatorTok{,}
        \VariableTok{player}
    \OperatorTok{)}
        \ControlFlowTok{return} \VariableTok{stack}\OperatorTok{:}\NormalTok{get\_count}\OperatorTok{()} \CommentTok{{-}{-} Hineinlegen erlauben}
    \KeywordTok{end}\OperatorTok{,}

    \VariableTok{allow\_take} \OperatorTok{=} \KeywordTok{function}\OperatorTok{(}
        \VariableTok{inv}\OperatorTok{,}
        \VariableTok{listname}\OperatorTok{,}
        \VariableTok{index}\OperatorTok{,}
        \VariableTok{stack}\OperatorTok{,}
        \VariableTok{player}
    \OperatorTok{)}
        \ControlFlowTok{return} \DecValTok{0} \CommentTok{{-}{-} Herausnehmen verbieten}
    \KeywordTok{end}\OperatorTok{,}

    \VariableTok{on\_put} \OperatorTok{=} \KeywordTok{function}\OperatorTok{(}\VariableTok{inv}\OperatorTok{,} \VariableTok{listname}\OperatorTok{,} \VariableTok{index}\OperatorTok{,} \VariableTok{stack}\OperatorTok{,} \VariableTok{player}\OperatorTok{)}
        \VariableTok{core}\OperatorTok{.}\NormalTok{chat\_send\_all}\OperatorTok{(}
            \VariableTok{player}\OperatorTok{:}\NormalTok{get\_player\_name}\OperatorTok{()}
                \OperatorTok{..} \StringTok{" hat "}
                \OperatorTok{..} \VariableTok{stack}\OperatorTok{:}\NormalTok{to\_string}\OperatorTok{()}
                \OperatorTok{..} \StringTok{" in die Spendenkiste gelegt, von "}
                \OperatorTok{..} \VariableTok{core}\OperatorTok{.}\NormalTok{pos\_to\_string}\OperatorTok{(}\VariableTok{player}\OperatorTok{:}\NormalTok{get\_pos}\OperatorTok{())}
        \OperatorTok{)}
    \KeywordTok{end}\OperatorTok{,}
\OperatorTok{\})}
\end{Highlighting}
\end{Shaded}

Berechtigungs-Callbacks -- also die mit \texttt{allow\_\allowbreak{}} am
Anfang -- geben die Anzahl der zu übertragenden Items zurück; 0
verhindert die Übertragung ganz.

Aktions-Callbacks -- die mit \texttt{on\_\allowbreak{}} am Anfang --
haben dagegen keinen Rückgabewert.

\section{Listen}\label{listen}

Mit Inventarlisten lassen sich mehrere Gitter an einem einzigen Ort
speichern. Das ist vor allem für Spieler nützlich, denn es gibt einige
Listen, die alle Spiele haben, etwa das Hauptinventar (\texttt{main})
und die Craftingfelder (\texttt{craft}).

\subsection{Größe und Breite}\label{gruxf6uxdfe-und-breite}

Listen haben eine Größe, also die Gesamtzahl der Felder, und eine
Breite, die nur intern in der Engine verwendet wird. Beim Anzeigen des
Inventars in einem Fenster spielt die Breite der Liste keine Rolle, weil
der Code des Fensters die Breite selbst festlegt.

\begin{Shaded}
\begin{Highlighting}[]
\ControlFlowTok{if} \VariableTok{inv}\OperatorTok{:}\NormalTok{set\_size}\OperatorTok{(}\StringTok{"main"}\OperatorTok{,} \DecValTok{32}\OperatorTok{)} \ControlFlowTok{then}
    \VariableTok{inv}\OperatorTok{:}\NormalTok{set\_width}\OperatorTok{(}\StringTok{"main"}\OperatorTok{,} \DecValTok{8}\OperatorTok{)}
    \FunctionTok{print}\OperatorTok{(}\StringTok{"Größe:  "} \OperatorTok{..} \VariableTok{inv}\OperatorTok{:}\NormalTok{get\_size}\OperatorTok{(}\StringTok{"main"}\OperatorTok{))}
    \FunctionTok{print}\OperatorTok{(}\StringTok{"Breite: "} \OperatorTok{..} \VariableTok{inv}\OperatorTok{:}\NormalTok{get\_width}\OperatorTok{(}\StringTok{"main"}\OperatorTok{))}
\ControlFlowTok{else}
    \FunctionTok{print}\OperatorTok{(}\StringTok{"Fehler: ungültige Liste oder Größe"}\OperatorTok{)}
\ControlFlowTok{end}
\end{Highlighting}
\end{Shaded}

\texttt{set\_\allowbreak{}size} schlägt fehl und gibt \texttt{false}
zurück, wenn der Listenname oder die Größe ungültig ist -- etwa wenn die
neue Größe zu klein für die Items ist, die gerade im Inventar liegen.

\subsection{Inhalt prüfen}\label{inhalt-pruxfcfen}

Mit \texttt{is\_\allowbreak{}empty} prüfst du, ob eine Liste überhaupt
Items enthält:

\begin{Shaded}
\begin{Highlighting}[]
\ControlFlowTok{if} \VariableTok{inv}\OperatorTok{:}\NormalTok{is\_empty}\OperatorTok{(}\StringTok{"main"}\OperatorTok{)} \ControlFlowTok{then}
    \FunctionTok{print}\OperatorTok{(}\StringTok{"Die Liste ist leer!"}\OperatorTok{)}
\ControlFlowTok{end}
\end{Highlighting}
\end{Shaded}

Mit \texttt{contains\_\allowbreak{}item} prüfst du, ob eine Liste ein
bestimmtes Item enthält:

\begin{Shaded}
\begin{Highlighting}[]
\ControlFlowTok{if} \VariableTok{inv}\OperatorTok{:}\NormalTok{contains\_item}\OperatorTok{(}\StringTok{"main"}\OperatorTok{,} \StringTok{"default:stone"}\OperatorTok{)} \ControlFlowTok{then}
    \FunctionTok{print}\OperatorTok{(}\StringTok{"Ich habe Stein gefunden!"}\OperatorTok{)}
\ControlFlowTok{end}
\end{Highlighting}
\end{Shaded}

\section{Inventare und ItemStacks
verändern}\label{inventare-und-itemstacks-veruxe4ndern}

\subsection{Zu einer Liste
hinzufügen}\label{zu-einer-liste-hinzufuxfcgen}

\texttt{add\_\allowbreak{}item} fügt Items zu einer Liste hinzu (hier
\texttt{\textquotedbl{}main\textquotedbl{}}). Im folgenden Beispiel wird
dabei auch die maximale Stapelgröße beachtet:

\begin{Shaded}
\begin{Highlighting}[]
\KeywordTok{local} \VariableTok{stack} \OperatorTok{=}\NormalTok{ ItemStack}\OperatorTok{(}\StringTok{"default:stone 99"}\OperatorTok{)}
\KeywordTok{local} \VariableTok{leftover} \OperatorTok{=} \VariableTok{inv}\OperatorTok{:}\NormalTok{add\_item}\OperatorTok{(}\StringTok{"main"}\OperatorTok{,} \VariableTok{stack}\OperatorTok{)}
\ControlFlowTok{if} \VariableTok{leftover}\OperatorTok{:}\NormalTok{get\_count}\OperatorTok{()} \OperatorTok{\textgreater{}} \DecValTok{0} \ControlFlowTok{then}
    \FunctionTok{print}\OperatorTok{(}
        \StringTok{"Inventar ist voll! "}
            \OperatorTok{..} \VariableTok{leftover}\OperatorTok{:}\NormalTok{get\_count}\OperatorTok{()}
            \OperatorTok{..} \StringTok{" Items wurden nicht hinzugefügt"}
    \OperatorTok{)}
\ControlFlowTok{end}
\end{Highlighting}
\end{Shaded}

\subsection{Items herausnehmen}\label{items-herausnehmen}

So entfernst du Items aus einer Liste:

\begin{Shaded}
\begin{Highlighting}[]
\KeywordTok{local} \VariableTok{taken} \OperatorTok{=} \VariableTok{inv}\OperatorTok{:}\NormalTok{remove\_item}\OperatorTok{(}\StringTok{"main"}\OperatorTok{,} \VariableTok{stack}\OperatorTok{)}
\FunctionTok{print}\OperatorTok{(}\StringTok{"Herausgenommen: "} \OperatorTok{..} \VariableTok{taken}\OperatorTok{:}\NormalTok{get\_count}\OperatorTok{())}
\end{Highlighting}
\end{Shaded}

\subsection{Stacks verändern}\label{stacks-veruxe4ndern}

Einzelne Stacks veränderst du, indem du sie zuerst holst:

\begin{Shaded}
\begin{Highlighting}[]
\KeywordTok{local} \VariableTok{stack} \OperatorTok{=} \VariableTok{inv}\OperatorTok{:}\NormalTok{get\_stack}\OperatorTok{(}\VariableTok{listname}\OperatorTok{,} \DecValTok{1}\OperatorTok{)}
\end{Highlighting}
\end{Shaded}

Dann änderst du sie, indem du Eigenschaften setzt oder Methoden
verwendest, die die Stapelgröße beachten:

\begin{Shaded}
\begin{Highlighting}[]
\KeywordTok{local} \VariableTok{stack} \OperatorTok{=}\NormalTok{ ItemStack}\OperatorTok{(}\StringTok{"default:stone 50"}\OperatorTok{)}
\KeywordTok{local} \VariableTok{to\_add} \OperatorTok{=}\NormalTok{ ItemStack}\OperatorTok{(}\StringTok{"default:stone 100"}\OperatorTok{)}
\KeywordTok{local} \VariableTok{leftover} \OperatorTok{=} \VariableTok{stack}\OperatorTok{:}\NormalTok{add\_item}\OperatorTok{(}\VariableTok{to\_add}\OperatorTok{)}
\KeywordTok{local} \VariableTok{taken} \OperatorTok{=} \VariableTok{stack}\OperatorTok{:}\NormalTok{take\_item}\OperatorTok{(}\DecValTok{19}\OperatorTok{)}

\FunctionTok{print}\OperatorTok{(}
    \StringTok{"Konnte "}
        \OperatorTok{..} \VariableTok{leftover}\OperatorTok{:}\NormalTok{get\_count}\OperatorTok{()}
        \OperatorTok{..} \StringTok{" Items nicht hinzufügen."}
\OperatorTok{)}
\CommentTok{{-}{-} \^{} ergibt 51}

\FunctionTok{print}\OperatorTok{(}\StringTok{"Habe "} \OperatorTok{..} \VariableTok{stack}\OperatorTok{:}\NormalTok{get\_count}\OperatorTok{()} \OperatorTok{..} \StringTok{" Items"}\OperatorTok{)}
\CommentTok{{-}{-} \^{} ergibt 80}
\CommentTok{{-}{-}   min(50+100, stack\_max) {-} 19 = 80}
\CommentTok{{-}{-}     mit stack\_max = 99}
\end{Highlighting}
\end{Shaded}

\texttt{add\_\allowbreak{}item} fügt einem ItemStack Items hinzu und
gibt alles zurück, was nicht hineingepasst hat.
\texttt{take\_\allowbreak{}item} nimmt höchstens die angegebene Anzahl
heraus, eventuell auch weniger, und gibt den herausgenommenen Stack
zurück.

Zum Schluss setzt du den ItemStack wieder ein:

\begin{Shaded}
\begin{Highlighting}[]
\VariableTok{inv}\OperatorTok{:}\NormalTok{set\_stack}\OperatorTok{(}\VariableTok{listname}\OperatorTok{,} \DecValTok{1}\OperatorTok{,} \VariableTok{stack}\OperatorTok{)}
\end{Highlighting}
\end{Shaded}

Die Felder einer Liste werden, wie in Lua üblich, ab 1 gezählt.

\section{Abnutzung}\label{abnutzung}

Werkzeuge können sich abnutzen. Die Abnutzung\index{abnutzung@Abnutzung}
wird als Balken angezeigt, und ist das Werkzeug vollständig abgenutzt,
zerbricht es. Die Abnutzung ist eine Zahl bis 65535; je höher sie ist,
desto stärker ist das Werkzeug abgenutzt.

Du veränderst sie mit \texttt{add\_\allowbreak{}wear()},
\texttt{get\_\allowbreak{}wear()} und
\texttt{set\_\allowbreak{}wear(wear)}.

\begin{Shaded}
\begin{Highlighting}[]
\KeywordTok{local} \VariableTok{stack} \OperatorTok{=}\NormalTok{ ItemStack}\OperatorTok{(}\StringTok{"default:pick\_mese"}\OperatorTok{)}
\KeywordTok{local} \VariableTok{max\_uses} \OperatorTok{=} \DecValTok{10}

\CommentTok{{-}{-} Passiert automatisch beim Abbauen mit einem Werkzeug:}
\CommentTok{{-}{-} erhöht die Abnutzung um eine Benutzung.}
\VariableTok{stack}\OperatorTok{:}\NormalTok{add\_wear}\OperatorTok{(}\DecValTok{65535} \OperatorTok{/} \OperatorTok{(}\VariableTok{max\_uses} \OperatorTok{{-}} \DecValTok{1}\OperatorTok{))}
\end{Highlighting}
\end{Shaded}

Beim Abbauen kann die Abnutzung eines Werkzeugs vom abgebauten Node
abhängen; \texttt{max\_\allowbreak{}uses} ist also je nach Node
verschieden.

\section{Lua-Tabellen}\label{lua-tabellen}

ItemStacks und Inventare lassen sich in Tabellen umwandeln und
umgekehrt. Das ist praktisch zum Kopieren und für Massenoperationen.

\begin{Shaded}
\begin{Highlighting}[]
\CommentTok{{-}{-} Ganzes Inventar}
\KeywordTok{local} \VariableTok{data} \OperatorTok{=} \VariableTok{inv1}\OperatorTok{:}\NormalTok{get\_lists}\OperatorTok{()}
\VariableTok{inv2}\OperatorTok{:}\NormalTok{set\_lists}\OperatorTok{(}\VariableTok{data}\OperatorTok{)}

\CommentTok{{-}{-} Eine Liste}
\KeywordTok{local} \VariableTok{listdata} \OperatorTok{=} \VariableTok{inv1}\OperatorTok{:}\NormalTok{get\_list}\OperatorTok{(}\StringTok{"main"}\OperatorTok{)}
\VariableTok{inv2}\OperatorTok{:}\NormalTok{set\_list}\OperatorTok{(}\StringTok{"main"}\OperatorTok{,} \VariableTok{listdata}\OperatorTok{)}
\end{Highlighting}
\end{Shaded}

Die von \texttt{get\_\allowbreak{}lists()} zurückgegebene Tabelle sieht
so aus:

\begin{Shaded}
\begin{Highlighting}[]
\OperatorTok{\{}
    \VariableTok{list\_one} \OperatorTok{=} \OperatorTok{\{}
        \VariableTok{ItemStack}\OperatorTok{,}
        \VariableTok{ItemStack}\OperatorTok{,}
        \VariableTok{ItemStack}\OperatorTok{,}
        \VariableTok{ItemStack}\OperatorTok{,}
        \CommentTok{{-}{-} inv:get\_size("list\_one") Elemente}
    \OperatorTok{\},}
    \VariableTok{list\_two} \OperatorTok{=} \OperatorTok{\{}
        \VariableTok{ItemStack}\OperatorTok{,}
        \VariableTok{ItemStack}\OperatorTok{,}
        \VariableTok{ItemStack}\OperatorTok{,}
        \VariableTok{ItemStack}\OperatorTok{,}
        \CommentTok{{-}{-} inv:get\_size("list\_two") Elemente}
    \OperatorTok{\},}
\OperatorTok{\},}
\end{Highlighting}
\end{Shaded}

\texttt{get\_\allowbreak{}list()} gibt eine einzelne Liste zurück, also
einfach eine Liste von ItemStacks.

Wichtig: Die \texttt{set}-Methoden oben verändern die Größe der Listen
nicht. Du kannst eine Liste also leeren, indem du sie auf eine leere
Tabelle setzt, ohne dass sie kleiner wird:

\begin{Shaded}
\begin{Highlighting}[]
\VariableTok{inv}\OperatorTok{:}\NormalTok{set\_list}\OperatorTok{(}\StringTok{"main"}\OperatorTok{,} \OperatorTok{\{\})}
\end{Highlighting}
\end{Shaded}
